\section*{Abstract}

Accessing government services in Bangladesh is often difficult due to confusing procedures, lack of clear and centralized information and reliance on middlemen. This prevents many citizens from accessing these services efficiently, resulting in unnecessary time delays and financial burdens. This thesis aims to solve this problem by developing an AI-Powered interactive system that provides citizen centric guidance for certain government services which will act as a digital bridge between citizens and government services. This system will inform users about required procedures and credentials of the respective services by offering dialogue and question-answering sessions. The outcome is a user-friendly, accessible system that empowers citizens with correct information, reduces dependency on intermediaries and simplifies interactions with government offices.

\vspace{1cm}
\textbf{Keywords:} Artificial Intelligence (AI), Government Services, Intermediaries, AI-Powered interactive system, Citizen centric guidance, Digital bridge, Dialogue and question-answering sessions, User-friendly and accessible system.
\pagebreak